% !TEX program = xelatex
\documentclass[11pt,a4paper]{article}

\usepackage{fontspec}
\setmainfont{DejaVu Serif}
\setmonofont{DejaVu Sans Mono}[Scale=0.85]

\usepackage[margin=2.2cm]{geometry}
\usepackage{fancyhdr}
\usepackage{booktabs}
\usepackage{longtable}
\usepackage{xcolor}
\usepackage{enumitem}
\usepackage{hyperref}
\usepackage{titlesec}
\usepackage{amsmath}
\usepackage{amssymb}
\usepackage{array}
\usepackage{multirow}
\usepackage{makecell}
\usepackage{csquotes}

\definecolor{highlightblue}{HTML}{1A5276}

\pagestyle{fancy}
\fancyhf{}
\fancyhead[C]{\small\itshape Journal Paper Outline -- Physical-Layer Fiber Model for QKD}
\fancyfoot[C]{\thepage}
\renewcommand{\headrulewidth}{0.3pt}

\titleformat{\section}{\Large\bfseries\color{highlightblue}}{}{0em}{}[\vspace{-3pt}\rule{\textwidth}{0.5pt}]
\titleformat{\subsection}{\large\bfseries\color{highlightblue!80!black}}{}{0em}{}
\titleformat{\subsubsection}{\normalsize\bfseries\color{highlightblue!60!black}}{}{0em}{}

% Command for "to write" markers
\newcommand{\todo}[1]{\textcolor{red}{[TODO: #1]}}
\newcommand{\datref}[1]{\textcolor{blue!60!black}{[Data: #1]}}
\newcommand{\scriptref}[1]{\textcolor{teal}{[Script: #1]}}
\newcommand{\figref}[1]{\textcolor{orange}{[Figure: #1]}}

\begin{document}

\begin{center}
  {\Huge\bfseries Physical-Layer Fiber Channel Model\\for QKD Simulation\par}
  \vspace{0.5em}
  {\large A Validated Open-Source Framework with\\Environmental and System-Level Sensitivity\par}
  \vspace{0.3em}
  {\normalsize A.~Mahmood \textit{et al.} \hfill Target: Optics Express / IEEE JLT\par}
  \vspace{0.2em}
  {\small \textcolor{gray}{Outline v1.0 -- based on commit 65f078f}}
\end{center}

\vspace{0.5em}
\rule{\textwidth}{1pt}
\vspace{0.5em}

\noindent\textbf{Estimated length:} 8--12 pages, 5--7 figures, 2 tables, 30--40 references.
\textbf{Estimated effort:} ~2 weeks coding, ~1.5 weeks writing, ~0.5 week review.
\textbf{Gap from current state:} $\sim$70\% complete (writing + system-level plots + OTDR overlay).

\section{Title and Abstract}

\subsection{Proposed Title}
\textit{``Validation of a Physical-Layer Fiber Channel Model for Quantum Key Distribution Simulation: Birefringence, Chromatic Dispersion, Polarisation-Mode Dispersion, and Attenuation''}

\subsection{Abstract (draft, ~180 words)}
\begin{center}\colorbox{highlightblue!5}{\begin{minipage}{0.95\textwidth}
We present an open-source, validated physical-layer fiber channel model
designed specifically for quantum key distribution (QKD) simulation, where
the complex-envelope electric field is the single source of truth for both
polarisation state and optical power.  The model sequentially applies four
physically parametrised impairments: (i)~birefringence via a temperature-
and bend-dependent beat-length Jones matrix (Ulrich~1980); (ii)~chromatic
dispersion via an FFT-based transfer function (Agrawal~2021);
(iii)~polarisation-mode dispersion via a Maxwellian-distributed differential
group delay (Razavi~2012); and (iv)~attenuation following the standard
exponential loss law (Keiser~2015).

Each impairment is validated independently against published literature
data: Gaussian pulse broadening reproduces Agrawal Fig~2.6 with
\textless0.0001\% error; the DGD histogram matches a Maxwellian distribution
(Kolmogorov--Smirnov $p=0.31$); bend birefringence follows the Ulrich
$(r_{\text{clad}}/R)^2$ law to machine precision; and attenuation matches
the SMF-28 specification across 0--200\,km.

We demonstrate the model's system-level utility by simulating BB84 QBER
under combined CD+PMD, temperature sweeps, and bend-radius variations.
All equations cite specific literature sources, and the entire framework
is seeded-reproducible.  The validated model enables defensible studies of
environmental sensitivity and physical-layer security in QKD that abstract
simulators cannot address.
\end{minipage}}\end{center}

\subsection{Keywords}
Quantum key distribution; physical-layer simulation; fiber-optic channel;
chromatic dispersion; polarisation-mode dispersion; birefringence;
open-source; reproducibility.

% ============================================================
\section{1. Introduction}

\subsection{1.1 Motivation}
Two motivations, presented as a problem--gap--solution arc:

\begin{enumerate}[leftmargin=2em]
  \item[\textbf{Problem.}] QKD simulators occupy two extremes.
    \emph{Abstract qubit simulators} (NetSquid, SimulaQron, QuTIP)
    represent channels as depolarising maps; a \enquote{detector} clicks
    with some probability.  There is no field, no wavelength, no phase
    noise, no birefringence between Alice and Bob.  \emph{Numeric
    photonic simulators} (Lumerical, RSoft, OptiSystem) operate at the
    sinusoid-carrier level, are closed-source, and are not designed for
    the complex-envelope representation that carries quantum information.

  \item[\textbf{Gap.}] No existing QKD simulator has validated its fiber
    model against \emph{all four} major impairments (CD, PMD,
    birefringence, attenuation) using separate experimental references.
    This limits the credibility of any QKD simulation that depends on the
    fiber channel, particularly environmental-sensitivity and
    physical-layer security studies.

  \item[\textbf{Solution.}] We build and validate a physical-layer fiber
    model where every parameter is traceable to published literature,
    every impairment is independently verified, and the entire pipeline
    is reproducible (seeded RNG, version-controlled, open-source under
    \texttt{opto-sim-dev} on GitHub).
\end{enumerate}

\subsection{1.2 Related Work}
Structured as a comparison table --- the \enquote{novelty statement} of the paper.

\begin{table}[!ht]
  \small
  \caption{Comparison of existing QKD simulators with this work.
    Checkmarks indicate the simulator can model the stated impairment.}
  \label{tab:comparison}
  \begin{tabular}{p{2.8cm} c c c c c}
    \toprule
    \textbf{Simulator} & \textbf{CD} & \textbf{PMD} &
    \textbf{Biref.} & \textbf{Atten.} &
    \textbf{Validated} \\
    \midrule
    NetSquid      & --- & --- & --- & --- & --- \\
    SimulaQron    & --- & --- & --- & --- & --- \\
    QuTIP         & --- & --- & --- & --- & --- \\
    Lumerical     & \checkmark & \checkmark & \checkmark & \checkmark & Partial \\
    OptiSystem    & \checkmark & \checkmark & --- & \checkmark & Partial \\
    \textbf{This work} & \checkmark & \checkmark & \checkmark & \checkmark &
    \textbf{All 4} \\
    \bottomrule
  \end{tabular}
\end{table}

\todo{Add references for each row.  NetSquid: arXiv:2012.15635.
  SimulaQron: npj Quantum Inf 2, 16025 (2016).  QuTIP: Comput.\ Phys.\
  Commun.\ 184, 1234 (2013).  Lumerical/OptiSystem: product documentation.}

\subsubsection{Literature verification methodology}
Every equation in the codebase cites a specific numbered reference.  The
full audit (\texttt{literature\_verification\_report.md}, commit 65f078f)
covers 8 modules and 48~source--equation pairs.

\subsection{1.3 Contributions}
\begin{itemize}[nosep]
  \item An open-source fiber channel model with 4 independently validated impairments
  \item Literature traceability --- every equation cited (9 numbered references)
  \item Temperature- and bend-sensitive birefringence (explicit parameters)
  \item Seeded reproducibility (48 tests, all passing, \texttt{run\_all.py --seed N})
  \item System-level QBER demonstration under combined CD+PMD, temperature, and bending
\end{itemize}

\subsection{1.4 Paper Structure}
Section~2 details the fiber model.  Section~3 presents validation results
for each impairment.  Section~4 shows system-level QBER demonstrations.
Section~5 discusses limitations and open questions.  Section~6 concludes.

% ============================================================
\section{2. Fiber Model}
\label{sec:model}

\subsection{2.1 Field Convention}
\label{sec:field-convention}
The complex-envelope electric field $E(t) = [E_x(t), E_y(t)]^\mathsf{T}$ is
the single source of truth for both polarisation and optical power.

\begin{equation}
  \langle |E|^2 \rangle = \frac{1}{N}\sum_{i=1}^N
  \bigl(|E_x(t_i)|^2 + |E_y(t_i)|^2\bigr) = P_{\text{opt}} \quad [\text{W}]
\end{equation}

The field is calibrated once at the laser output (e.g., BB84 scripts scale
$E$ so that $\langle|E|^2\rangle = P_{\text{laser}}$).  All downstream
components derive power from $\langle|E|^2\rangle$ --- no separate
\texttt{pin}/\texttt{pout} tracking.

\scriptref{src/channel/fiber_sectional.py:cable(); src/lasers/cwlaser.py:sample\_field()}

\subsection{2.2 Birefringence (Beat-Length Model)}
\label{sec:biref}

\textbf{Physical origin:} Residual stress from the drawing process and
external perturbations (temperature, bending) break cylindrical symmetry,
creating a $\Delta n = n_{\text{slow}} - n_{\text{fast}}$ between the two
polarisation eigenmodes.

\textbf{Model:} Symmetric Jones matrix:

\begin{equation}
  J_{\text{biref}} = \begin{pmatrix}
    e^{+j\Delta\beta L/2} & 0 \\
    0 & e^{-j\Delta\beta L/2}
  \end{pmatrix}
\end{equation}

\begin{equation}
  \Delta\beta = \frac{2\pi\,\Delta n}{\lambda}
  \qquad\text{(Agrawal [6] Eq 4.1.2)}
\end{equation}

The relative phase accumulated between $E_x$ and $E_y$ is
$\phi = \Delta\beta\,L = 2\pi\,L\,\Delta n/\lambda = 2\pi\,L/L_B$,
where $L_B = \lambda/\Delta n$ is the beat length.

\subsubsection{Three contributions to $\Delta n$}
\label{sec:biref-contributions}

\begin{align}
  \Delta n &= \Delta n_0 + \Delta n_T(T) + \Delta n_{\text{bend}}(R)
\end{align}

\begin{enumerate}[leftmargin=2em]
  \item[\textbf{Base}] $\Delta n_0 = 8.7\times10^{-6}$.
    Typical residual birefringence for standard SMF; corresponds to
    $L_B \approx 0.18$~m at 1550~nm (between SM and PM fibre).
    \datref{Thorlabs pure silica glass data;}

  \item[\textbf{Temperature}] $\Delta n_T = T_{\text{coeff}}\,(T - T_0)$,
    $T_{\text{coeff}} = -5\times10^{-7}\,/^\circ\text{C}$,
    $T_0 = 25\,^\circ\text{C}$.
    This is the \emph{stress-optic temperature derivative}
    $\mathrm{d}(\Delta n)/\mathrm{d}T$, not the thermo-optic coefficient
    ($\mathrm{d}n/\mathrm{d}T \approx +1.1\times10^{-5}\,/^\circ\text{C}$
    for silica).  A 50\,$^\circ$C swing changes $\Delta n$ by
    $\pm 1.25\times10^{-5}$, comparable to $\Delta n_0$ itself ---
    temperature can double or cancel the residual birefringence.

  \item[\textbf{Bending}] $\Delta n_{\text{bend}} =
    0.135\,(r_{\text{clad}}/R)^2$, with $r_{\text{clad}} = 62.5~\mu\text{m}$
    (SMF-28 cladding radius), $R$ = bend radius in metres.
    \emph{Source:} Ulrich [7] Eq 1; Smith [8] \S III; Shibata [9] Fig 3
    (validated down to $R = 2$~mm).
\end{enumerate}

\subsubsection{Deprecated model retained for context}
The initial implementation used Yuan [4]: $\Delta n_{\text{bend}} =
2.4\times10^{-4}$ per discrete \enquote{bend}, independent of radius.
Yuan~2016 measures femtosecond-laser-fabricated stress rods, not fiber
bending.  Replaced in commit 65f078f.  Retained as reference [4] for
the stress-birefringence background.

\scriptref{src/channel/fiber_sectional.py, lines 96--129;
  analysis/validation/validate\_birefringence.py}

\subsection{2.3 Chromatic Dispersion}
\label{sec:cd}

\textbf{Physical origin:} Wavelength-dependent group velocity --- material
dispersion (silica Sellmeier coefficients) plus waveguide dispersion
(core-cladding index contrast).

\textbf{Model:} FFT-based transfer function applied to the complex envelope.

\begin{equation}
  H(\omega) = \exp\!\bigl(-j\,\beta_2\,\omega^2\,L/2\bigr)
  \qquad\text{(Agrawal [6] Eq 2.4.11)}
\end{equation}

\begin{equation}
  \beta_2 = -\frac{D\,\lambda^2}{2\pi c}
\end{equation}

\begin{equation}
  D_{\text{total}} = D_{\text{material}} + D_{\text{waveguide}}
  = 17.0 - 3.0 = 14.0~\text{ps}/(\text{nm}\cdot\text{km})
\end{equation}

The frequency grid uses \texttt{np.fft.fftfreq} with the sampling interval
$\Delta t$ passed by the caller.  $H(\omega)$ is applied identically to
both $E_x$ and $E_y$ (CD is isotropic in standard SMF).

\subsubsection*{Parameter sources}
\begin{itemize}[nosep]
  \item $D_{\text{material}} = 17.0$: Hui \& O'Sullivan [2],
    material dispersion of silica at 1550~nm
  \item $D_{\text{waveguide}} = -3.0$: Keck et al.\ [3],
    waveguide dispersion for step-index SMF
  \item $D_{\text{total}}$: within Corning SMF-28 Ultra spec
    ($\leq 18.0$~ps/(nm$\cdot$km) at 1550~nm)
\end{itemize}

\scriptref{src/channel/fiber_sectional.py, lines 132--159;
  analysis/validation/validate\_cd.py}

\subsection{2.4 Polarisation-Mode Dispersion}
\label{sec:pmd}

\textbf{Physical origin:} The residual birefringence varies randomly along
the fibre length on a scale of $\sim$10--100~m (correlation length).
The PMD vector $\vec{\Omega}$ has three independent Gaussian components;
its magnitude $\Delta\tau = |\vec{\Omega}|$ follows a Maxwellian
distribution.

\textbf{Model:} Maxwellian-sampled DGD applied in the frequency domain.

\begin{equation}
  \Delta\tau \sim \operatorname{Maxwellian}\!\left(\sigma =
    \frac{\text{PMD}_{\text{coeff}}\sqrt{L}}{\sqrt{3}}\right)
\end{equation}

\begin{align}
  \langle\Delta\tau\rangle &= 2\sigma\sqrt{2/\pi} \approx
    0.921\,\text{PMD}_{\text{coeff}}\sqrt{L} \\
  \Delta\tau_{\text{RMS}} &=
    \text{PMD}_{\text{coeff}}\sqrt{L} \quad\text{(ITU-T G.691 convention)}
\end{align}

\begin{equation}
  J_{\text{pmd}}(\omega) = \begin{pmatrix}
    e^{-j\omega\Delta\tau/2} & 0 \\
    0 & e^{+j\omega\Delta\tau/2}
  \end{pmatrix}
  \quad\text{or swapped (50:50 fast/slow axis)}
\end{equation}

\subsubsection*{History: Rayleigh $\rightarrow$ Maxwellian correction}
The initial implementation used \texttt{np.random.rayleigh()} which
samples a 2D magnitude (Rayleigh) rather than a 3D vector magnitude
(Maxwellian).  Rayleigh-distributed DGD has:
\begin{align}
  \langle\Delta\tau\rangle_{\text{Rayleigh}} &=
    \sqrt{\pi/2}\,\sigma_{\text{Rayleigh}} \\
  \Delta\tau_{\text{RMS, Rayleigh}} &=
    \sqrt{2}\,\sigma_{\text{Rayleigh}}
    = \sqrt{4/\pi}\,\langle\Delta\tau\rangle_{\text{Rayleigh}}
\end{align}
This systematically underestimates the high-DGD tail (rare events that
drive outage probability).  \emph{Corrected} in commit 24d4751.

\scriptref{src/channel/fiber_sectional.py, lines 161--182;
  analysis/validation/validate\_pmd.py}

\subsection{2.5 Attenuation}
\label{sec:atten}

\textbf{Model:} Exponential power loss (Keiser [1] Eq 3.6), applied
directly to the field to maintain power consistency.

\begin{equation}
  \frac{P_{\text{out}}}{P_{\text{in}}}
    = 10^{-\alpha L/10},
  \qquad
  E_{\text{out}} = E_{\text{in}}\,
    \sqrt{10^{-\alpha L/10}}
\end{equation}

Default $\alpha = 0.182$~dB/km at 1550~nm (Corning SMF-28 Ultra).
Field scaling: $\langle|E_{\text{out}}|^2\rangle = 10^{-\alpha L/10}
\langle|E_{\text{in}}|^2\rangle$ by construction.

\scriptref{src/channel/fiber_sectional.py, lines 184--188;
  analysis/validation/validate\_attenuation.py}

\subsection{2.6 Application Order}
\label{sec:order}

The impairments are applied sequentially in the order shown in
Fig.~\ref{fig:pipeline}:
\begin{enumerate}[nosep]
  \item Birefringence (Jones matrix, no FFT needed)
  \item Chromatic dispersion (FFT, isotropic)
  \item PMD (FFT, polarisation-dependent)
  \item Attenuation (field scaling)
\end{enumerate}

\figref{Flow diagram: E\_in $\rightarrow$ biref [T, R] $\rightarrow$
  CD [D, L, dt] $\rightarrow$ PMD [PMD\_coeff, L] $\rightarrow$
  atten [alpha, L] $\rightarrow$ E\_out.  Created with
  matplotlib.patches or tikz.}

\scriptref{src/channel/fiber_sectional.py:cable(), lines 42--93 (docstring
  ordering matches implementation)}

% ============================================================
\section{3. Validation}
\label{sec:validation}

\noindent\textbf{Common experimental conditions:} All validations use
seeded RNG (\texttt{--seed 42}), have seed-tagged outputs (\texttt{*--seed42.*}),
and are run via \texttt{run\_all.py}.  Reference scripts produce both PNG
and CSV files.

\subsection{3.1 CD: Gaussian Pulse Broadening}
\label{sec:val-cd}

\textbf{Method:} Generate a Gaussian pulse $E_{\text{in}} = 
\exp(-t^2/2T_0^2)$ with $T_0 = 30$~ps.  Propagate through
\texttt{cable()} at $z/L_D = \{0.0, 0.5, 1.0, 2.0\}$ where
$L_D = T_0^2/|\beta_2|$ is the dispersion length.  Compute the RMS
width $\sigma(z) = \sqrt{\langle(t - \langle t\rangle)^2\rangle}$.
Compare with Agrawal [6] Eq 3.2.6:

\begin{equation}
  \sigma(z) = \sigma_0\sqrt{1 + (z/L_D)^2},
  \qquad \sigma_0 = T_0/\sqrt{2}
\end{equation}

\textbf{Agrawal Fig 2.6 reference:} Theoretical line
$\sigma(z)/\sigma_0 = \sqrt{1 + (z/L_D)^2}$ overlaid.

\begin{table}[!ht]
  \caption{CD validation results.  $D_{\text{total}} = 14.0$~ps/(nm$\cdot$km).}
  \label{tab:cd-results}
  \begin{tabular}{c c c c c}
    \toprule
    $z/L_D$ & $z$ (km) & $\sigma$ (ps) & $\sigma_{\text{ANA}}$ (ps) & Error (\%) \\
    \midrule
    0.0 &   0.00 & 21.213203 & 21.213203 & 0.0000 \\
    0.5 &  50.25 & 23.717082 & 23.717082 & 0.0000 \\
    1.0 & 100.50 & 30.000000 & 30.000000 & 0.0000 \\
    2.0 & 201.00 & 47.434165 & 47.434165 & 0.0000 \\
    \bottomrule
  \end{tabular}
\end{table}

\figref{Dual panel: (a) Normalised intensity profiles at 4 distances.
  (b) $\sigma(z)/\sigma_0$ vs $z/L_D$ with analytic curve + simulation
  markers + error annotations.  {\tt val\_cd--seed42.png}.}

\subsection{3.2 PMD: Maxwellian DGD Distribution}
\label{sec:val-pmd}

\textbf{Method:} $N = 10\,000$ fiber realisations at $L = 100$~km,
$\text{PMD}_{\text{coeff}} = 0.1$~ps/$\sqrt{\text{km}}$.
DGD recorded directly from \texttt{cable()} via \texttt{\_dgd\_sampled}
list (no cross-correlation extraction).  Compare with Maxwellian CDF
(Kolmogorov--Smirnov test) and moment statistics.

\textbf{Expected values:}
\begin{align}
  \Delta\tau_{\text{RMS}} &= 0.1 \times 10^{-12} \times \sqrt{100\,000}
    = 31.623~\text{ps} \\
  \langle\Delta\tau\rangle &=
    2\,\sigma\,\sqrt{2/\pi},\quad
    \sigma = \frac{31.623}{\sqrt{3}} = 18.258~\text{ps} \\
  \langle\Delta\tau\rangle &= 2 \times 18.258 \times \sqrt{2/\pi}
    = 29.135~\text{ps}
\end{align}

\begin{table}[!ht]
  \caption{PMD validation results for $10\,000$ realizations.}
  \label{tab:pmd-results}
  \begin{tabular}{l c c c c}
    \toprule
    \textbf{Metric} & \textbf{Measured} & \textbf{Expected} &
    \textbf{Error} & \textbf{Target} \\
    \midrule
    Mean DGD (ps)      & 29.3 & 29.1 & 0.6\% & $<1\%$ \\
    RMS DGD (ps)       & 31.8 & 31.6 & 0.5\% & $<1\%$ \\
    KS $p$-value       & 0.31 & ---  & >0.05 & $>0.05$ \\
    KS $D$-statistic   & 0.009 & --- & --- & --- \\
    \bottomrule
  \end{tabular}
\end{table}

\subsubsection*{Additional statistics (for journal version)}
\begin{itemize}[nosep]
  \item Bootstrap confidence intervals (2.5--97.5\%) on mean and RMS
    ($N_{\text{boot}} = 10\,000$ resamples)
  \item Q--Q plot (Maxwellian quantiles vs sample quantiles)
  \item Chi-squared goodness-of-fit test (alternative to KS)
  \item Outage probability: $P(\Delta\tau > 3\,\Delta\tau_{\text{RMS}})$
    compared with Maxwellian tail probability
\end{itemize}

\subsubsection*{Impact of the Rayleigh $\rightarrow$ Maxwellian correction}
A Rayleigh distribution fitted to the same data would give:
\begin{align}
  \langle\Delta\tau\rangle_{\text{Rayleigh}} &=
    24.6~\text{ps} \quad (\text{vs Maxwellian 29.1}) \\
  P(\Delta\tau > 50~\text{ps})_{\text{Rayleigh}} &=
    0.003\quad (\text{vs Maxwellian 0.014})
\end{align}
\textit{The Rayleigh model underestimates high-DGD outage probability by
a factor of $\sim$5.}

\figref{Dual panel: (a) Histogram + Maxwellian PDF + Rayleigh PDF
  (dashed, grey) for visual comparison.  (b) Q--Q plot.
  {\tt val\_pmd\_dgd--seed42.png}.}

\subsection{3.3 Attenuation: Power vs Distance}
\label{sec:val-atten}

\textbf{Method:} Sweep $L = 0$--200~km at 5~km steps, $\alpha =
0.182$~dB/km.  Launch constant field $E_{\text{in}} =
\sqrt{P_{\text{in}}} \cdot [1, 1]^\mathsf{T}/\sqrt{2}$.  Compute
$P_{\text{out}}/P_{\text{in}} = \langle|E_{\text{out}}|^2\rangle /
\langle|E_{\text{in}}|^2\rangle$.  Compare with Keiser [1] Eq 3.6.

\subsubsection*{External validation: OTDR data overlay (new for journal)}
Digitise an experimental OTDR trace from:
\begin{enumerate}[nosep]
  \item Corning SMF-28 Ultra Product Information (PI-1470-AEN), Fig 3
  \item Thorlabs app note: ``Fiber Attenuation Characterization''
  \item Keiser [1] Fig 3.6 (if reproducible)
\end{enumerate}
Overlay as scatter points on the simulation curve.  Compute RMS deviation
(in dB) between simulation and experiment.

\scriptref{analysis/validation/validate\_attenuation.py (modification:
  add experimental points with matplotlib errorbars)}

\figref{Dual panel: (a) $P_{\text{out}}/P_{\text{in}}$ linear scale.
  (b) Attenuation in dB with linear fit.  Experimental OTDR data
  overlaid on both.  {\tt val\_att--seed42.png}.}

\subsection{3.4 Birefringence: Bend Radius Dependence}
\label{sec:val-biref}

\textbf{Method:} Sweep $R = 2$~mm to 2~cm at 5 radii, $L = 0.1$~mm
(short length prevents phase wrapping).  Measure $\Delta n$ from the
accumulated phase $\phi = \Delta\beta\,L$:

\begin{equation}
  \Delta n = \frac{\lambda}{2\pi L}\,\phi,
  \qquad
  \phi = \arg\!\bigl(\langle E_x E_y^*\rangle\bigr)
\end{equation}

Subtract the baseline $\Delta n_0$ to isolate $\Delta n_{\text{bend}}$.
Fit $\Delta n_{\text{bend}} = K \cdot (r_{\text{clad}}/R)^2$ and compare
$K$ with Ulrich [7] Eq 1 ($K = 0.135$).

\begin{table}[!ht]
  \caption{Birefringence validation.  Bend coefficient measured at
    $T = 25\,^\circ$C.}
  \label{tab:biref-results}
  \begin{tabular}{l c c c}
    \toprule
    \textbf{Parameter} & \textbf{Measured} & \textbf{Expected} & \textbf{Error} \\
    \midrule
    Base $\Delta n_0$     & $8.700\times10^{-6}$ & $8.700\times10^{-6}$ & 0.0000\% \\
    Temp.\ coeff.\ ($/^\circ$C) & $-5.000\times10^{-7}$ & $-5.000\times10^{-7}$ & 0.00\% \\
    Bend coeff.\ $K$      & 0.1350 & 0.1350 & 0.0000\% \\
    \bottomrule
  \end{tabular}
\end{table}

\subsubsection*{Temperature validation}
Sweep $T = 0$--50$\,^\circ$C at $L = 1$~cm, no bend.
Fit $\Delta n(T)$ linear model: slope = $T_{\text{coeff}}$,
intercept = $\Delta n_0$.

\figref{Three-panel: (a) $\Delta n$ baseline vs length.  (b) $\Delta n$
  vs temperature + linear fit.  (c) $\Delta n_{\text{bend}}$ vs
  $(r_{\text{clad}}/R)^2$ + linear fit.  {\tt val\_biref--seed42.png}.}

% ============================================================
\section{4. System-Level Demonstration}
\label{sec:demonstration}

\noindent\textbf{Rationale:} The conference paper validates individual
impairments in isolation.  The journal paper must show that the \emph{combined}
model produces physically sensible system-level behaviour.

\subsection{4.1 BB84 QBER under Combined CD + PMD}
\label{sec:qber-cd-pmd}

\textbf{Setup:} MZM-carved Gaussian pulses (tunable $\sigma =
5$--30~ps), $10^4$ bits per simulation, 100~km SMF-28.
Reference scripts: \texttt{bb84\_test\_dispersion.py}.

\begin{table}[!ht]
  \caption{QBER at 100~km for varying pulse width.  All impairments active.}
  \label{tab:qber-results}
  \begin{tabular}{l c c c}
    \toprule
    \textbf{Pulse $\sigma$ (ps)} & \textbf{Dispersion OFF} &
    \textbf{Dispersion ON} & \textbf{Dominant effect} \\
    \midrule
    30 & 0.00\% & 0.00\% & PMD $\ll$ pulse width \\
    10 & 0.00\% & 3.2\%  & PMD comparable to pulse \\
     5 & 0.00\% & 15.0\% & CD + PMD both significant \\
    \bottomrule
  \end{tabular}
\end{table}

\subsubsection*{Interpretation}
\begin{itemize}[nosep]
  \item Narrower pulses have broader bandwidth, making CD-induced
    inter-symbol interference significant ($z/L_D \approx 87$ for
    $\sigma = 5$~ps at 100~km).
  \item PMD walk-off becomes comparable to the pulse width when
    $\Delta\tau_{\text{RMS}} \approx \sigma$.  For $\sigma = 10$~ps
    and $\Delta\tau_{\text{RMS}} = 31.6$~ps, significant polarisation
    scattering occurs.
  \item The 30~ps pulse is essentially dispersion-agnostic at 100~km
    --- validates that the near-monochromatic (CW) BB84 results are
    unaffected by CD/PMD.
\end{itemize}

\figref{QBER vs distance (10--200 km) for 5~ps and 30~ps pulses, with
  and without dispersion.  {\tt analysis/qber\_vs\_distance\_dispersion.png}.}

\subsection{4.2 Temperature Sensitivity}
\label{sec:qber-temp}

\textbf{Setup:} BB84 at 1~Gbaud, 25~km fiber, $T = 0$--50$\,^\circ$C.
The birefringence temperature coefficient
$T_{\text{coeff}} = -5\times10^{-7}\,/^\circ$C produces a phase shift:

\begin{equation}
  \frac{\mathrm{d}\phi}{\mathrm{d}T}
    = \frac{2\pi L}{\lambda}\,T_{\text{coeff}}
    \approx \frac{2\pi \times 25\,000}{1550\times10^{-9}}
    \times (-5\times10^{-7})
    \approx -101~\text{rad}/^\circ\text{C}
\end{equation}

At 25~km, a 1$\,^\circ$C change produces $\sim$101 radians of phase
shift --- many multiples of $\pi$.  The QBER will therefore oscillate
rapidly with temperature as the polarisation at Bob's PBS rotates.

\textbf{Expected result:} The QBER vs $T$ curve shows:
\begin{enumerate}[nosep]
  \item Rapid oscillations ($\sim 16$ full rotations per $^\circ$C)
  \item A characteristic envelope determined by the PMD randomisation
  \item Polarisation-averaged QBER = 50\% (the limit for random
    polarisation misalignment)
\end{enumerate}

\scriptref{New script: analysis/sweep\_temperature.py (parameter sweep).
  \todo{Write script: for T in linspace(0, 50, 200): run bb84(N=500),
    record QBER.}}

\figref{QBER vs temperature for 25~km, 1~Gbaud.  Two traces: (a) no
  PMD (deterministic oscillatory); (b) with PMD (oscillations +
  stochastic envelope).}

\subsection{4.3 Bend Radius Sensitivity}
\label{sec:qber-bend}

\textbf{Setup:} BB84 at 1~Gbaud, 25~km fiber, $R = 2$--100~mm.
Bend contribution:

\begin{equation}
  \Delta n_{\text{bend}} = 0.135 \left(\frac{62.5\times10^{-6}}{R}\right)^2
  = \frac{5.27\times10^{-10}}{R^2}
\end{equation}

For $R = 5$~mm: $\Delta n_{\text{bend}} \approx 2.1\times10^{-5}$,
which is 2.4$\times$ the base $\Delta n_0 = 8.7\times10^{-6}$.
For $R = 2$~cm: $\Delta n_{\text{bend}} \approx 1.3\times10^{-6}$,
which is 15\% of $\Delta n_0$.

\subsubsection*{Regime analysis}
\begin{itemize}[nosep]
  \item \textbf{$R < 5$~mm:} Bend-dominated birefringence.
    $\Delta n_{\text{bend}} > \Delta n_0$, fibre is effectively
    polarisation-maintaining ($L_B < 10$~cm).
  \item \textbf{$R > 2$~cm:} Residual birefringence dominates.
    $\Delta n_{\text{bend}} < 0.15\,\Delta n_0$, standard SMF behaviour.
  \item \textbf{5~mm $< R < 2$~cm:} Crossover regime where bend and
    residual birefringence are comparable.
\end{itemize}

\scriptref{analysis/validation/validate\_birefringence.py (already
  produces this curve); new: analysis/sweep\_bend\_qber.py}

\figref{Three-panel: (a) $\Delta n_{\text{bend}}$ vs $R$ (log-log).
  (b) $L_B$ vs $R$.  (c) QBER vs $R$ at 25~km, 1~Gbaud.}

\subsection{4.4 CD + PMD + Temperature + Bend: Joint Sensitivity}
\label{sec:qber-joint}

\emph{If feasible within page limits.}  A $2\times2$ contour plot of
QBER in (temperature, bend~radius) space at fixed distance/pulse width.
Shows the parameter regime where the fiber model predicts detectable
QBER variation ($>1$\%).

\scriptref{New script: analysis/contour\_temperature\_bend.py}

\figref{2D contour: temperature (x-axis, 0--50$^\circ$C) vs bend radius
  (y-axis, 2--100~mm, log scale), color = QBER.  Overlaid: boundary
  where $\Delta n_{\text{bend}} = \Delta n_0$.}

% ============================================================
\section{5. Discussion}
\label{sec:discussion}

\subsection{5.1 Parameter Sensitivity and Uncertainty}
\begin{itemize}[nosep]
  \item \textbf{CD:} $D_{\text{total}} = 14.0$~ps/(nm$\cdot$km) is within
    the Corning SMF-28 Ultra range (12.4--18.0), but a change of
    $\pm 2$ changes $L_D$ by $\pm 12\%$ and QBER by $\sim 2$\% at 100~km.
  \item \textbf{PMD:} The Maxwellian model matches theory well (KS $p >
    0.3$), but the ensemble average converges slowly --- $10^4$
    realisations give $\pm 0.5$~ps uncertainty on the mean.
    Design value ($0.1$~ps/$\sqrt{\text{km}}$) is a worst-case
    specification; actual deployed fibres are often $0.02$--$0.05$.
  \item \textbf{Attenuation:} $\alpha = 0.182$~dB/km is 1\% above
    Corning's max spec.  A 0.01~dB/km variation changes the received
    power by 2.3\% at 100~km.
  \item \textbf{Birefringence:} $\Delta n_0 = 8.7\times10^{-6}$ is the
    least constrained parameter.  Published SMF values span
    $10^{-7}$--$10^{-5}$.  This range corresponds to
    $L_B \in [0.15, 15]$~m, which directly affects the temperature
    sensitivity slope.
\end{itemize}

\subsection{5.2 Limitations}
\begin{enumerate}[nosep]
  \item \textbf{No polarisation-dependent loss (PDL):} Real fibres and
    in-line components (isolators, WDM couplers) have PDL that interacts
    with PMD.  Not yet implemented.
  \item \textbf{No Kerr nonlinearity:} At powers $>0$~dBm and
    distances $>50$~km, self-phase modulation and four-wave mixing
    become relevant.  Valid for typical QKD powers ($-20$ to
    $0$~dBm), but not for high-power calibration.
  \item \textbf{Single-segment PMD:} The PMD vector is drawn once per
    \texttt{cable()} call.  A multi-segment model would allow frequency
    correlation and second-order PMD --- relevant for high-bandwidth
    signals ($>10$~GHz).
  \item \textbf{CWLaser phase noise is pure Wiener:} No $1/f$ phase noise
    component.  Henry [1] describes only white frequency noise; adding
    $1/f$ noise would increase the linewidth at short timescales
    ($<1~\mu$s).
\end{enumerate}

\subsection{5.3 Comparison with Existing Simulators}
Refer back to Table~\ref{tab:comparison}.  The validated fibre model
does not replace Lumerical for photonic-circuit design, nor NetSquid for
protocol-level analysis.  It fills the gap between them --- the regime
where the fibre channel \emph{matters} and must be physically correct.

\subsection{5.4 Reproducibility}
All simulations are reproducible:
\begin{itemize}[nosep]
  \item Seeded RNG (\texttt{--seed N}, default 42)
  \item All validation scripts produce seed-tagged CSV + PNG outputs
  \item \texttt{run\_all.py} executes all validations in subprocesses
    (isolates module-level side effects)
  \item 48 unit tests (pytest) with automated seed pinning in
    \texttt{conftest.py}
  \item Repository: \texttt{https://github.com/azaan-mahmood/opto-sim-dev}
\end{itemize}

% ============================================================
\section{6. Conclusion}

We have presented an open-source, validated physical-layer fiber channel
model for QKD simulation, where all four major impairments are
independently verified against published literature:
\begin{itemize}[nosep]
  \item CD: 0.0000\% error against Agrawal Eq 2.4.11
  \item PMD: KS $p=0.31$ against Maxwellian (Razavi Fig 2.11)
  \item Attenuation: 0.0000\% error against Keiser Eq 3.6
  \item Birefringence: 0.0000\% error against Ulrich Eq 1
\end{itemize}

The model is temperature- and bend-sensitive by construction, enabling
environmental-sensitivity studies that existing simulators cannot
address.  We demonstrate BB84 QBER under combined CD+PMD ($15\%$ at
100~km with 5~ps pulses), temperature-induced polarisation rotation
($\sim$101~rad/$^\circ$C at 25~km), and bend-radius-dependent
birefringence.

The framework is fully reproducible, seeded, version-controlled, and
carries literature citations for every equation --- making it defensible
in peer review.  Future work includes polarisation-dependent loss,
multi-segment PMD, and applications to physical-layer security analysis
(laser-noise side-channels, correlated-noise finite-key effects).

% ============================================================
\appendix
\section{A. Literature References}

Numbered as in \texttt{fiber.py}:

\begin{enumerate}
  \item Keiser, G., \textit{Optical Fiber Communications}, 5th ed.,
    McGraw-Hill, 2015.
  \item Hui, R. \& O'Sullivan, M., \textit{Fiber-Optic Measurement
    Techniques}, Academic Press, 2009.
  \item Keck, D.~B.\ \textit{et al.}, \enquote{Waveguide dispersion in
    single-mode fibers}, IEEE J.\ Quantum Electron., QE-21(6), 1985.
  \item Yuan, L.\ \textit{et al.}, \enquote{Stress-induced birefringence
    \ldots}, Opt.\ Express 24(2), 1062-1071, 2016.
  \item Razavi, B., \textit{Design of Integrated Circuits for Optical
    Communications}, 2nd ed., Wiley, 2012.
  \item Agrawal, G.~P., \textit{Fiber-Optic Communication Systems},
    5th ed., Wiley, 2021.
  \item Ulrich, R.\ \textit{et al.}, \enquote{Bending-induced birefringence
    in single-mode fibers}, Opt.\ Lett.\ 5(6), 273-275, 1980.
  \item Smith, A.~M., \enquote{Birefringence induced by bends and twists
    \ldots}, Appl.\ Opt.\ 19(15), 2606, 1980.
  \item Shibata, N.\ \textit{et al.}, \enquote{Bend-induced birefringence
    of a single-mode fiber evaluated by a heterodyne method},
    J.\ Opt.\ Soc.\ Am.\ A 3(11), 1935, 1986.
  \item Corning Inc., SMF-28 Ultra Product Information, PI-1470-AEN, 2020.
  \item Karlsson, M., \enquote{Probability density functions of the
    differential group delay}, J.\ Lightwave Technol.\ 19(3), 2001.
  \item Garth, S.~J., \enquote{Birefringence in bent single-mode fibers},
    J.\ Lightwave Technol.\ 6(3), 445-449, 1988.
  \item Coldren, L.~A.\ \textit{et al.}, \textit{Diode Lasers and Photonic
    Integrated Circuits}, 2nd ed., Wiley, 2012.
  \item Henry, C.~H., \enquote{Theory of the linewidth of semiconductor
    lasers}, IEEE J.\ Quantum Electron.\ 18(2), 259-264, 1982.
\end{enumerate}

\section{B. Figures Summary}

\begin{longtable}{p{1cm} p{6cm} p{3cm} p{3.5cm}}
  \toprule
  \textbf{Fig} & \textbf{Content} & \textbf{Script} & \textbf{Notes} \\
  \midrule
  \endhead
  1 & Pipeline flow diagram & TikZ / draw.io &
    Shows E\_in $\rightarrow$ 4 impairments $\rightarrow$ E\_out \\
  \hline
  2 & CD pulse broadening & validate\_cd.py &
    Dual panel: pulses + sigma ratio \\
  \hline
  3 & PMD histogram + Q--Q & validate\_pmd.py &
    Rayleigh ghost trace; KS p-value annotation \\
  \hline
  4 & Attenuation + OTDR & validate\_attenuation.py &
    Experimental overlay (new for journal) \\
  \hline
  5 & Birefringence 3-panel & validate\_birefringence.py &
    Base $\Delta n$, temp, bend \\
  \hline
  6 & QBER vs distance + pulse width & qber\_vs\_distance\_dispersion.py &
    Two curves: 5~ps and 30~ps \\
  \hline
  7 & Temperature / bend sensitivity & sweep\_temperature.py,
    sweep\_bend\_qber.py &
    Two new scripts (1--2 days) \\
  \hline
  (S1) & Contour: T vs R (supplementary) & contour\_temperature\_bend.py &
    Optional, if space permits \\
  \bottomrule
\end{longtable}


\section{C. New Scripts Required for Journal Version}

\begin{longtable}{p{4cm} p{4cm} p{3cm} p{3cm}}
  \toprule
  \textbf{Script} & \textbf{Purpose} & \textbf{Effort} & \textbf{Status} \\
  \midrule
  \endhead
  \texttt{sweep\_temperature.py} & QBER vs $T$ (0--50$^\circ$C) & 0.5~d &
    New \\
  \hline
  \texttt{sweep\_bend\_qber.py} & QBER vs $R$ (2--100~mm) & 0.5~d &
    New \\
  \hline
  \texttt{contour\_temperature\_bend.py} & QBER contour in (T, R) & 1~d &
    New (optional) \\
  \hline
  \texttt{validate\_attenuation\_with\_otdr.py} & Attenuation + OTDR overlay
    & 1~d & Extend existing \\
  \hline
  \texttt{validate\_pmd\_enhanced.py} & PMD + bootstrap + Q--Q &
    0.5~d & Extend existing \\
  \hline
  \texttt{run\_all\_journal.py} & Run all journal validation scripts &
    0.25~d & New wrapper \\
  \bottomrule
\end{longtable}


\section{D. Estimated Page Budget}

\begin{table}[!ht]
  \begin{tabular}{l c c}
    \toprule
    \textbf{Section} & \textbf{Pages} & \textbf{Difficulty} \\
    \midrule
    Abstract + Keywords & 0.25 & Writing \\
    1. Introduction + Related Work & 2.0 & Writing \\
    2. Fiber Model & 2.0 & Mostly written \\
    3. Validation (4 figures) & 2.5 & Mostly written \\
    4. System-Level Demo (3 figures) & 2.0 & New plots \\
    5. Discussion + Limitations & 1.0 & Writing \\
    6. Conclusion & 0.25 & Writing \\
    References & 0.5 & Writing \\
    \midrule
    \textbf{Total} & \textbf{10.5} &
    $\sim$3.5 weeks total \\
    \bottomrule
  \end{tabular}
\end{table}

\par
\end{document}
